\begin{thebibliography}{150}
\bibitem{ref1}
El Houssaini et al. (2022). Location-Aware IoT-Enabled WSNs for Landslide Early Warning. \textit{Electronics}. Available online: https://doi.org/10.3390/electronics11233971 (accessed on 17 October 2026).
\bibitem{ref2}
Giri et al. (2019). WSN System for Landslide Monitoring and Warning. \textit{IEEE Trans. Instrum. Meas.}. Available online: https://doi.org/10.1109/tim.2018.2861999 (accessed on 17 October 2026).
\bibitem{ref3}
Rani et al. (2025). IoT-Enabled Real-Time Landslide Monitoring using Waterbody Sensors. \textit{ICTMIM 2025}. Available online: https://doi.org/10.1109/ictmim65579.2025.10988011 (accessed on 17 October 2026).
\bibitem{ref4}
Gowtham et al. (2026). Low-Power LoRaWAN IoT Landslide EWS with Random Forest. \textit{ICAISS 2026}. Available online: https://doi.org/10.1109/icaiss68683.2026.11525861 (accessed on 17 October 2026).
\bibitem{ref5}
Amirthavarshini et al. (2025). Harnessing LoRa for Real-Time Landslide Monitoring in Kerala. \textit{Geoscience Frontiers}. Available online: https://doi.org/10.1016/j.gsf.2025.102242 (accessed on 17 October 2026).
\bibitem{ref6}
Souza et al. (2022). IoT-Based System for Landslides Warnings Using Real Time Slope Monitoring. \textit{COTB'22}. Available online: https://doi.org/10.14210/cotb.v13.p126-133 (accessed on 17 October 2026).
\bibitem{ref7}
D. J. et al. (2025). I-ELWS: IoT Driven Early Landslide Warning System. \textit{ETNCC 2025}. Available online: https://doi.org/10.1109/etncc66224.2025.11299651 (accessed on 17 October 2026).
\bibitem{ref8}
P. C. et al. (2021). Real-time landslide EWS through low cost soil and rainfall sensors. \textit{Materials Today: Proceedings}. Available online: https://doi.org/10.1016/j.matpr.2021.02.456 (accessed on 17 October 2026).
\bibitem{ref9}
Chaulya et al. (2024). Landslide monitoring and prediction system using geosensors and WSN. \textit{Discover Geoscience}. Available online: https://doi.org/10.1007/s44288-024-00007-3 (accessed on 17 October 2026).
\bibitem{ref10}
Susanto et al. (2024). Landslide Monitoring and Warning System Based on WSN. \textit{ELKHA}. Available online: https://doi.org/10.26418/elkha.v16i2.83489 (accessed on 17 October 2026).
\bibitem{ref11}
Thein et al. (2020). Real-Time Monitoring and EWS for Landslide Preventing in Myanmar. \textit{GCCE 2020}. Available online: https://doi.org/10.1109/gcce50665.2020.9291809 (accessed on 17 October 2026).
\bibitem{ref12}
Ramesh (2009). Real-Time WSN for Landslide Detection. \textit{SensorComm 2009}. Available online: https://doi.org/10.1109/sensorcomm.2009.67 (accessed on 17 October 2026).
\bibitem{ref13}
Devi (2025). Developing a Low-Cost Landslide EWS using IoT and Sensor Networks. \textit{ICSCSS 2025}. Available online: https://doi.org/10.1109/icscss64956.2025.11501071 (accessed on 17 October 2026).
\bibitem{ref14}
Gamperl et al. (2021). IoT Geosensor Network for Cost-Effective LEWS (Inform@Risk, Colombia). \textit{Sensors}. Available online: https://doi.org/10.3390/s21082609 (accessed on 17 October 2026).
\bibitem{ref15}
Mahfudh et al. (2026). Low-Cost IoT-Based Landslide EWS. \textit{CENTIVE Proceedings}. Available online: https://doi.org/10.20895/centive.v2025i1.562 (accessed on 17 October 2026).
\bibitem{ref16}
Rawat et al. (2024). LANDSLIDE MONITOR: a real-time landslide monitoring system. \textit{Environ. Earth Sci.}. Available online: https://doi.org/10.1007/s12665-024-11526-0 (accessed on 17 October 2026).
\bibitem{ref17}
Galley et al. (2017). WSN Based Landslide Detection and EWS. \textit{Zorig Melong}. Available online: https://doi.org/10.17102/zmv3.i1.006 (accessed on 17 October 2026).
\bibitem{ref18}
Abraham et al. (2020). IoT-Based Geotechnical Monitoring, Darjeeling Himalayas. \textit{Sensors}. Available online: https://doi.org/10.3390/s20092611 (accessed on 17 October 2026).
\bibitem{ref19}
Das et al. (2020). Real-Time Threshold-based Landslide Prediction using WSN. \textit{ICCE-Taiwan 2020}. Available online: https://doi.org/10.1109/icce-taiwan49838.2020.9258181 (accessed on 17 October 2026).
\bibitem{ref20}
Ramesh et al. (2017). WSNs for Early Warning of Landslides: Decade Long Deployment. Available online: https://doi.org/10.1007/978-3-319-53487-9\_4 (accessed on 17 October 2026).
\bibitem{ref21}
Ali et al. (2012). Development of WSN for slope monitoring (Hulu Langat, MALAYSIA, 2011 landslide, 16 fatalities). \textit{ICIAS 2012}. Available online: https://doi.org/10.1109/icias.2012.6306203 (accessed on 17 October 2026).
\bibitem{ref22}
Ragnoli et al. (2023). WSN-Based Rockfall and Landslide Monitoring Systems: A Review. \textit{Sensors}. Available online: https://doi.org/10.3390/s23167278 (accessed on 17 October 2026).
\bibitem{ref23}
Bagwari et al. (2021). Low-Cost Sensor and LoRaWAN Opportunities for Landslide Monitoring: A Review. \textit{IEEE Access}. Available online: https://doi.org/10.1109/access.2021.3137841 (accessed on 17 October 2026).
\bibitem{ref24}
Jeong et al. (2019). Effectiveness of a WSN System for Landslide Monitoring (Seoul). \textit{IEEE Access}. Available online: https://doi.org/10.1109/access.2019.2960570 (accessed on 17 October 2026).
\bibitem{ref25}
Marino et al. (2023). IoT-Based Low-Cost Sensor Network for Hydrological Monitoring of Landslide-Prone Areas. \textit{Sensors}. Available online: https://doi.org/10.3390/s23042299 (accessed on 17 October 2026).
\bibitem{ref26}
Debnath et al. (2025). Smart IoT-Based Distributed System for Remote Detection of Slope Movements (highway embankments). \textit{ICNC 2025}. Available online: https://doi.org/10.1109/icnc64010.2025.10993553 (accessed on 17 October 2026).
\bibitem{ref27}
Ahmed et al. (2016). Real-time monitoring of partially stable slopes using WSN. \textit{IC-GET 2016}. Available online: https://doi.org/10.1109/get.2016.7916638 (accessed on 17 October 2026).
\bibitem{ref28}
Vu et al. (2018). WSN in landslide monitoring system with remote data management. \textit{Measurement}. Available online: https://doi.org/10.1016/j.measurement.2018.01.002 (accessed on 17 October 2026).
\bibitem{ref29}
Kumar M. et al. (2022). Accurate IoT Based Slope Instability Sensing System for Landslide Detection. \textit{IEEE Sensors J.}. Available online: https://doi.org/10.1109/jsen.2022.3189903 (accessed on 17 October 2026).
\bibitem{ref30}
Joshi et al. (2023). WSN-Based Smart Landslide Monitoring Device (SMARTMODE). \textit{IEEE Trans. Instrum. Meas.}. Available online: https://doi.org/10.1109/tim.2023.3269746 (accessed on 17 October 2026).
\bibitem{ref31}
Idris et al. (n.d.). Smart Monitoring and Warning Landslide System using IoT (Bukit Bauk Dungun, MALAYSIA).
\bibitem{ref32}
Piciullo et al. (2017). Threshold-based regional EW model for rainfall-induced landslides. \textit{Landslides}. Available online: https://doi.org/10.1007/s10346-016-0750-2 (accessed on 17 October 2026).
\bibitem{ref33}
Liu et al. (2024). Regional EW model based on slope unit, Chongqing. \textit{Engineering Geology}. Available online: https://doi.org/10.1016/j.enggeo.2024.107464 (accessed on 17 October 2026).
\bibitem{ref34}
Dou et al. (2024). Regional dynamic EW model, Fujian. \textit{Geomatics, Nat. Hazards and Risk}. Available online: https://doi.org/10.1080/19475705.2024.2435510 (accessed on 17 October 2026).
\bibitem{ref35}
Guzzetti et al. (2024). Landslide predictions through combined rainfall threshold models. \textit{Landslides}. Available online: https://doi.org/10.1007/s10346-024-02340-7 (accessed on 17 October 2026).
\bibitem{ref36}
Huang et al. (2022). Regional rainfall-induced landslide hazard warning based on susceptibility mapping. \textit{Geomorphology}. Available online: https://doi.org/10.1016/j.geomorph.2022.108236 (accessed on 17 October 2026).
\bibitem{ref37}
Zhang et al. (2026). Dynamic EW Model Coupling Slope Units and Rainfall Thresholds, Jinyang. \textit{Geosciences}. Available online: https://doi.org/10.3390/geosciences16090367 (accessed on 17 October 2026).
\bibitem{ref38}
Cao et al. (2022). Hourly shallow landslide warning model combining susceptibility and rainfall threshold. \textit{J. Mountain Sci.}. Available online: https://doi.org/10.1007/s11629-022-7370-1 (accessed on 17 October 2026).
\bibitem{ref39}
Song et al. (2025). Refined meteorological EW based on temporal probability model. \textit{Chinese J. Rock Mech. Eng.}. Available online: https://doi.org/10.3724/1000-6915.jrme.2024.0805 (accessed on 17 October 2026).
\bibitem{ref40}
Wang et al. (2026). Dynamic EW model based on refined threshold analysis, Tanjiawan. \textit{PLOS One}. Available online: https://doi.org/10.1371/journal.pone.0339689 (accessed on 17 October 2026).
\bibitem{ref41}
Zhan et al. (2022). EW model and model test verification of rainfall-induced shallow landslide. \textit{Bull. Eng. Geol. Environ.}. Available online: https://doi.org/10.1007/s10064-022-02827-4 (accessed on 17 October 2026).
\bibitem{ref42}
Nocentini et al. (2023). Optimization of rainfall thresholds through false alarm reduction. \textit{Landslides}. Available online: https://doi.org/10.1007/s10346-023-02176-7 (accessed on 17 October 2026).
\bibitem{ref43}
Sun et al. (2024). Probabilistic thresholds for regional rainfall induced landslides. \textit{Computers and Geotechnics}. Available online: https://doi.org/10.1016/j.compgeo.2023.106040 (accessed on 17 October 2026).
\bibitem{ref44}
Naidu et al. (2017). EWS for shallow landslides using rainfall threshold and slope stability analysis. \textit{Geoscience Frontiers}. Available online: https://doi.org/10.1016/j.gsf.2017.10.008 (accessed on 17 October 2026).
\bibitem{ref45}
Dou et al. (2023). Rainfall EW threshold and spatial distribution in China. \textit{Rock Mechanics Bulletin}. Available online: https://doi.org/10.1016/j.rockmb.2023.100056 (accessed on 17 October 2026).
\bibitem{ref46}
Dixit et al. (2024). Numerical-model-derived I-D thresholds for debris flows, Himalayan catchment. \textit{NHESS}. Available online: https://doi.org/10.5194/nhess-24-465-2024 (accessed on 17 October 2026).
\bibitem{ref47}
Lee et al. (2021). Optimal rainfall thresholds for LEWS, Chuncheon, Korea. \textit{Landslides}. Available online: https://doi.org/10.1007/s10346-020-01603-3 (accessed on 17 October 2026).
\bibitem{ref48}
Wang et al. (2026). Physical-based rainfall landslide prediction model from slope-unit I-D curves. \textit{NHESS}. Available online: https://doi.org/10.5194/nhess-26-2367-2026 (accessed on 17 October 2026).
\bibitem{ref49}
Sarkar et al. (2023). Event-duration rainfall threshold model, Uttarkashi, Himalayas. \textit{Int. J. Earth Sci.}. Available online: https://doi.org/10.1007/s00531-023-02337-y (accessed on 17 October 2026).
\bibitem{ref50}
Li et al. (2026). Sustained Self-Powered Real-Time Vibration Monitoring via Nonlinear Harvesting. \textit{Advanced Science}. Available online: https://doi.org/10.1002/advs.76967 (accessed on 17 October 2026).
\bibitem{ref51}
Rubeš et al. (2021). Self-powered wireless vibration sensor node using vibration energy harvester. \textit{Mech. Syst. Signal Process.}. Available online: https://doi.org/10.1016/j.ymssp.2021.107890 (accessed on 17 October 2026).
\bibitem{ref52}
Izhar et al. (2023). Hybrid acoustic, vibration, and wind energy harvester for self-powered WSN. \textit{Energy Conversion and Management}. Available online: https://doi.org/10.1016/j.enconman.2022.116635 (accessed on 17 October 2026).
\bibitem{ref53}
Wang et al. (2024). Sensor-enabled piezoelectric geocable–geogrid for slope rehabilitation and monitoring. \textit{Engineering Geology}. Available online: https://doi.org/10.1016/j.enggeo.2024.107780 (accessed on 17 October 2026).
\bibitem{ref54}
Wang et al. (2026). Reinforcement and microseismic monitoring of slopes with self-sensing smart geosynthetic. \textit{Geosynthetics International}. Available online: https://doi.org/10.1680/jgein.25.00212 (accessed on 17 October 2026).
\bibitem{ref55}
Rigo et al. (2024). Piezoelectric Sensors as Energy Harvesters for Ultra Low-Power IoT. \textit{Sensors}. Available online: https://doi.org/10.3390/s24082587 (accessed on 17 October 2026).
\bibitem{ref56}
Wang et al. (2024). Self-sustained/self-wakeup wireless vibration sensors, hybrid harvesting. \textit{Applied Energy}. Available online: https://doi.org/10.1016/j.apenergy.2023.122207 (accessed on 17 October 2026).
\bibitem{ref57}
Yang et al. (2025). Hybrid multiple-layer piezoelectric-solar harvester for self-powered sensing. \textit{Sustainable Energy Technologies and Assessments}. Available online: https://doi.org/10.1016/j.seta.2024.104157 (accessed on 17 October 2026).
\bibitem{ref58}
Li et al. (2021). Piezoelectric Energy Harvesting: Materials, Structures, Applications. \textit{Small Structures}. Available online: https://doi.org/10.1002/sstr.202100128 (accessed on 17 October 2026).
\bibitem{ref59}
Ng et al. (2005). Sensitivity Analysis and Energy Harvesting for a Self-Powered Piezoelectric Sensor. \textit{J. Intelligent Material Systems and Structures}. Available online: https://doi.org/10.1177/1045389x05053151 (accessed on 17 October 2026).
\bibitem{ref60}
Lo et al. (2022). Self-powered SECE piezoelectric harvesting from shock excitations. \textit{Mech. Syst. Signal Process.}. Available online: https://doi.org/10.1016/j.ymssp.2022.109123 (accessed on 17 October 2026).
\bibitem{ref61}
Biswal et al. (2024). Self-Powered Accelerometer Using Circular Hole Piezo Harvester. \textit{IEEE Sensors Letters}. Available online: https://doi.org/10.1109/lsens.2024.3485540 (accessed on 17 October 2026).
\bibitem{ref62}
Kuang et al. (2021). Strongly coupled piezoelectric harvesters for self-powered condition monitoring. \textit{Energy Conversion and Management}. Available online: https://doi.org/10.1016/j.enconman.2021.114129 (accessed on 17 October 2026).
\bibitem{ref63}
Zou (2023). Piezoelectric self-powered wireless sensor in vibration environment monitoring. \textit{Results in Engineering}. Available online: https://doi.org/10.1016/j.rineng.2023.101492 (accessed on 17 October 2026).
\bibitem{ref64}
Das Mahapatra et al. (2021). Piezoelectric Materials for Energy Harvesting and Sensing: Roadmap. \textit{Advanced Science}. Available online: https://doi.org/10.1002/advs.202100864 (accessed on 17 October 2026).
\bibitem{ref66}
Lepipas et al. (2024). Miniature water flow harvester, Savonius-type microturbine. \textit{Smart Materials and Structures}. Available online: https://doi.org/10.1088/1361-665x/ad1c3f (accessed on 17 October 2026).
\bibitem{ref67}
Guiffard et al. (2021). Self-powered wireless sensor with flexible aero-piezoelectric harvester. \textit{Renewable Energy}. Available online: https://doi.org/10.1016/j.renene.2021.11.113 (accessed on 17 October 2026).
\bibitem{ref68}
Zhang et al. (2024). Triboelectric-EM hybrid generator, Savonius flapping wing, low-velocity water flow. \textit{Applied Energy}. Available online: https://doi.org/10.1016/j.apenergy.2023.122512 (accessed on 17 October 2026).
\bibitem{ref69}
Rahman et al. (2025). Survey of flow-based energy harvesters for sustainable WSN nodes. \textit{J. Renewable and Sustainable Energy}. Available online: https://doi.org/10.1063/5.0237597 (accessed on 17 October 2026).
\bibitem{ref70}
Zhao et al. (2021). Hydrokinetic piezoelectric harvesting by wake-induced vibration. \textit{Energy}. Available online: https://doi.org/10.1016/j.energy.2020.119722 (accessed on 17 October 2026).
\bibitem{ref71}
Kamenar et al. (2016). Harvesting river flow energy for WSN technology (GOLDFISH project). \textit{Microsystem Technologies}. Available online: https://doi.org/10.1007/s00542-015-2778-y (accessed on 17 October 2026).
\bibitem{ref72}
Li et al. (2022). Self-powered wireless water-monitoring sensor, electromagnetic-pendulum harvester. \textit{Energy}. Available online: https://doi.org/10.1016/j.energy.2022.123883 (accessed on 17 October 2026).
\bibitem{ref73}
Vijayaregunathan et al. (2024). Harvesting hydraulic energy for self-sustainable IoT flow measurement. \textit{Energy Sources Part A}. Available online: https://doi.org/10.1080/15567036.2024.2334922 (accessed on 17 October 2026).
\bibitem{ref74}
Yunus et al. (2026). Design of a Wireless Sensor Node Powered by Kinetic Energy Harvesters. \textit{Applied Research}. Available online: https://doi.org/10.1002/appl.70124 (accessed on 17 October 2026).
\bibitem{ref75}
Zhou et al. (2026). Small-scale water energy harvesting for sustainably-powered distributed electronics. \textit{Communications Materials}. Available online: https://doi.org/10.1038/s43246-026-01137-6 (accessed on 17 October 2026).
\bibitem{ref76}
Sarker et al. (2024). Micro energy harvesting for IoT platform: Review. \textit{Heliyon}. Available online: https://doi.org/10.1016/j.heliyon.2024.e27778 (accessed on 17 October 2026).
\bibitem{ref77}
Hafizh et al. (2023). Vortex-induced-vibration self-tunable airfoil piezoelectric harvester for underwater sensing. \textit{Ocean Engineering}. Available online: https://doi.org/10.1016/j.oceaneng.2022.113467 (accessed on 17 October 2026).
\bibitem{ref78}
Azevedo et al. (2012). Energy harvesting from wind and water for autonomous WSN nodes. \textit{IET Circuits Devices Syst.}. Available online: https://doi.org/10.1049/iet-cds.2011.0287 (accessed on 17 October 2026).
\bibitem{ref79}
Williams et al. (2021). Survey of Energy Harvesting Technologies for WSNs. \textit{IEEE Access}. Available online: https://doi.org/10.1109/access.2021.3083697 (accessed on 17 October 2026).
\bibitem{ref81}
Ragnoli et al. (2020). Autonomous Low-Power LoRa-Based Flood-Monitoring System. \textit{J. Low Power Electronics and Applications}. Available online: https://doi.org/10.3390/jlpea10020015 (accessed on 17 October 2026).
\bibitem{ref82}
Tsakmakis et al. (2022). Adaptive LoRaWAN MAC Protocol for Event Detection Applications. \textit{Sensors}. Available online: https://doi.org/10.3390/s22093538 (accessed on 17 October 2026).
\bibitem{ref83}
Pan et al. (2021). UAV-Aided Emergency Environmental Monitoring, LoRa Mesh. \textit{IEEE IoT Journal}. Available online: https://doi.org/10.1109/jiot.2021.3095494 (accessed on 17 October 2026).
\bibitem{ref84}
Sendra et al. (2020). LoRaWAN Network for Fire Monitoring in Rural Environments. \textit{Electronics}. Available online: https://doi.org/10.3390/electronics9030531 (accessed on 17 October 2026).
\bibitem{ref86}
Ragnoli et al. (2022). LoRa-Based WSN for Rockfall and Landslide Monitoring, Pantelleria Island. \textit{J. Low Power Electronics and Applications}. Available online: https://doi.org/10.3390/jlpea12030047 (accessed on 17 October 2026).
\bibitem{ref87}
Wu et al. (2023). Low-Cost Low-Power LoRa Mesh Network for Large-Scale Environmental Sensing (CottonCandy). \textit{IEEE IoT Journal}. Available online: https://doi.org/10.1109/jiot.2023.3270237 (accessed on 17 October 2026).
\bibitem{ref88}
Finochietto et al. (2022). Reducing Operational Expenses of LoRaWAN-Based IoRT Applications. \textit{Sensors}. Available online: https://doi.org/10.3390/s22207778 (accessed on 17 October 2026).
\bibitem{ref89}
Bagwari et al. (2022). Disaster Monitoring based on IoT and Long Range Assisted Framework. \textit{J. Physics: Conf. Series}. Available online: https://doi.org/10.1088/1742-6596/2327/1/012020 (accessed on 17 October 2026).
\bibitem{ref90}
Ragnoli et al. (2022). Flood monitoring: LoRa-based case study, L'Aquila. \textit{PRIME 2022}. Available online: https://doi.org/10.1109/prime55000.2022.9816747 (accessed on 17 October 2026).
\bibitem{ref91}
Abdul Latiff (2019). Environmental Monitoring System Based on LoRa Technology in Island. \textit{ICSigSys 2019}. Available online: https://doi.org/10.1109/icsigsys.2019.8811066 (accessed on 17 October 2026).
\bibitem{ref92}
Chapungo et al. (2025). LoRaWAN Connectivity Reliability in Remote Rural Areas of Mozambique (source of the ~89\% PDR lossy-link finding used in Phase 3 synthesis). \textit{Sensors}. Available online: https://doi.org/10.3390/s25196027 (accessed on 17 October 2026).
\bibitem{ref93}
Zhao et al. (2023). Forest IoT Using LPWAN, Solar-Powered Micro-Weather Station. \textit{Sensors}. Available online: https://doi.org/10.3390/s23104809 (accessed on 17 October 2026).
\bibitem{ref94}
Tamang et al. (2022). Reliable Low-Latency LoRaWAN Solution for Environmental Monitoring at Major Accident Risk Sites. \textit{Sensors}. Available online: https://doi.org/10.3390/s22062372 (accessed on 17 October 2026).
\bibitem{ref95}
Liya et al. (2026). Energy efficient cyber-physical control of renewable microgrids, edge-AI, blockchain. \textit{Scientific Reports}. Available online: https://doi.org/10.1038/s41598-026-50169-y (accessed on 17 October 2026).
\bibitem{ref96}
Silva et al. (2023). IoT-based EMS for AC microgrids with grid and security constraints. \textit{Applied Energy}. Available online: https://doi.org/10.1016/j.apenergy.2023.120904 (accessed on 17 October 2026).
\bibitem{ref97}
Wu et al. (2024). Microgrids EM considering net-zero energy, IoT digital twin. \textit{Sustainable Energy Technologies and Assessments}. Available online: https://doi.org/10.1016/j.seta.2024.103621 (accessed on 17 October 2026).
\bibitem{ref98}
Cicceri et al. (2023). Deep Learning-Driven Self-Conscious Distributed CPS for Renewable Energy Communities. \textit{Sensors}. Available online: https://doi.org/10.3390/s23094549 (accessed on 17 October 2026).
\bibitem{ref99}
Asef et al. (2023). SIEMS: Secure Intelligent EMS for Industrial IoT Applications. \textit{IEEE Trans. Industrial Informatics}. Available online: https://doi.org/10.1109/tii.2022.3165890 (accessed on 17 October 2026).
\bibitem{ref100}
Tasmant et al. (2025). Review of ML and IoT-based EMS for AC microgrids. \textit{Comput. Electr. Eng.}. Available online: https://doi.org/10.1016/j.compeleceng.2025.110563 (accessed on 17 October 2026).
\bibitem{ref101}
Lee \& Wang (2024). A Cyber-Physical Testbed for IoT Microgrid Design and Validation [PRIMARY REFERENCE PAPER]. \textit{Electronics}. Available online: https://doi.org/10.3390/electronics13071181 (accessed on 17 October 2026).
\bibitem{ref102}
Albarakati et al. (2022). Microgrid energy management and monitoring systems: A comprehensive review. \textit{Frontiers in Energy Research}. Available online: https://doi.org/10.3389/fenrg.2022.1097858 (accessed on 17 October 2026).
\bibitem{ref103}
Xu et al. (2021). Blockchain-based secure energy policy and management of renewable smart microgrids. \textit{Sustainable Cities and Society}. Available online: https://doi.org/10.1016/j.scs.2021.103010 (accessed on 17 October 2026).
\bibitem{ref104}
Sitharthan et al. (2023). Smart microgrid with IoT for energy management (PV/Wind/Piezo hybrid). \textit{Comput. Electr. Eng.}. Available online: https://doi.org/10.1016/j.compeleceng.2022.108556 (accessed on 17 October 2026).
\bibitem{ref105}
Mallala et al. (2025). IoT-enabled community microgrid EM with battery degradation consideration. \textit{J. Energy Storage}. Available online: https://doi.org/10.1016/j.est.2024.114739 (accessed on 17 October 2026).
\bibitem{ref107}
Samanta et al. (2020). IoT-based smart EM in VRFB-integrated bio-solar microgrid. \textit{J. Energy Storage}. Available online: https://doi.org/10.1016/j.est.2020.101967 (accessed on 17 October 2026).
\bibitem{ref108}
Aljafari et al. (2022). Optimization of DC, AC, Hybrid AC/DC Microgrid-Based IoT Systems: A Review. \textit{Energies}. Available online: https://doi.org/10.3390/en15186813 (accessed on 17 October 2026).
\bibitem{ref109}
Sedhom et al. (2021). IoT-based optimal demand side management for smart microgrid. \textit{Int. J. Electrical Power \& Energy Systems}. Available online: https://doi.org/10.1016/j.ijepes.2020.106674 (accessed on 17 October 2026).
\bibitem{ref110}
Zhou et al. (2021). Multi-microgrid EMS: Architecture, Communication, and Scheduling Strategies. \textit{J. Modern Power Systems and Clean Energy}. Available online: https://doi.org/10.35833/mpce.2019.000237 (accessed on 17 October 2026).
\bibitem{ref111}
Ojo et al. (2025). Microgrids' Control Strategies and Real-Time Monitoring Systems: Comprehensive Review. \textit{Energies}. Available online: https://doi.org/10.3390/en18133576 (accessed on 17 October 2026).
\bibitem{ref112}
Alhasnawi et al. (2023). New communication platform for smart EMS using MILP. \textit{Energy Systems}. Available online: https://doi.org/10.1007/s12667-023-00591-2 (accessed on 17 October 2026).
\bibitem{ref113}
Bahbahani et al. (2017). Cooperative Clustering Protocol With Duty Cycling for EH-WSN. \textit{IEEE Trans. Wireless Communications}. Available online: https://doi.org/10.1109/twc.2017.2762674 (accessed on 17 October 2026).
\bibitem{ref114}
Shagari et al. (2020). Heterogeneous Energy and Traffic Aware Sleep-Awake Cluster-Based Routing. \textit{IEEE Access}. Available online: https://doi.org/10.1109/access.2020.2965206 (accessed on 17 October 2026).
\bibitem{ref115}
Ramkumar (2025). Adaptive Duty-Cycled Cooperative Clustering Model for EH-WSN. \textit{CNC 2025}. Available online: https://doi.org/10.1109/cnc68716.2025.11484711 (accessed on 17 October 2026).
\bibitem{ref116}
Liu et al. (2023). DCC-IACJS: bio-inspired duty cycle-based clustering for energy-efficient WSN. \textit{J. King Saud University CIS}. Available online: https://doi.org/10.1016/j.jksuci.2023.01.015 (accessed on 17 October 2026).
\bibitem{ref117}
V. et al. (2025). Energy-Aware MAC Protocols with Load Balancing for Solar-Powered WSN. \textit{ICDSNS 2025}. Available online: https://doi.org/10.1109/icdsns65743.2025.11168589 (accessed on 17 October 2026).
\bibitem{ref118}
Tolani et al. (2023). Energy-Efficient Duty-Cycle Hybrid MAC Protocol for WSN. \textit{Wireless Personal Communications}. Available online: https://doi.org/10.1007/s11277-023-10448-7 (accessed on 17 October 2026).
\bibitem{ref119}
Chauhan et al. (2022). Variable duty cycle aware energy efficient clustering strategy for WSN. \textit{J. Ambient Intelligence and Humanized Computing}. Available online: https://doi.org/10.1007/s12652-022-04363-1 (accessed on 17 October 2026).
\bibitem{ref120}
Gong et al. (2023). Centralized Energy-Efficient Clustering Protocol for WSN. \textit{IEEE Sensors J.}. Available online: https://doi.org/10.1109/jsen.2022.3224180 (accessed on 17 October 2026).
\bibitem{ref121}
Yang et al. (2025). DRL-Driven Workload Offloading for Energy-Efficient Solar-Powered WSN. \textit{ACM CoNEXT Workshop}. Available online: https://doi.org/10.1145/3769698.3771226 (accessed on 17 October 2026).
\bibitem{ref122}
Jalalinejad et al. (2024). Hybrid Multi-Hop Clustering and Energy-Aware Routing for EH-WSN. \textit{IEEE Access}. Available online: https://doi.org/10.1109/access.2024.3458795 (accessed on 17 October 2026).
\bibitem{ref123}
Zhao et al. (2023). Load-Aware Energy-Efficient Clustering Algorithm in Sensor-Cloud. \textit{J. Grid Computing}. Available online: https://doi.org/10.1007/s10723-023-09683-w (accessed on 17 October 2026).
\bibitem{ref124}
Hady (2020). Duty Cycling Centralized Hierarchical Routing Protocol for WSN. \textit{Comput. Syst. Sci. Eng.}. Available online: https://doi.org/10.32604/csse.2020.35.347 (accessed on 17 October 2026).
\bibitem{ref125}
Yadav et al. (2023). Clustering-Based Energy Efficient Task Offloading for Sustainable Fog Computing. \textit{IEEE Trans. Sustainable Computing}. Available online: https://doi.org/10.1109/tsusc.2022.3186585 (accessed on 17 October 2026).
\bibitem{ref126}
González et al. (2024). Evolutionary, Swarm, Iterative-Based Task-Offloading Optimization for WSN Battery Life. \textit{IEEE Sensors J.}. Available online: https://doi.org/10.1109/jsen.2024.3419558 (accessed on 17 October 2026).
\bibitem{ref127}
Li et al. (2020). Energy-Aware Task Offloading in the Internet of Things. \textit{IEEE Wireless Communications}. Available online: https://doi.org/10.1109/mwc.001.1900495 (accessed on 17 October 2026).
\bibitem{ref128}
Zhang et al. (2018). Peer-to-Peer energy trading in a Microgrid. \textit{Applied Energy}. Available online: https://doi.org/10.1016/j.apenergy.2018.03.010 (accessed on 17 October 2026).
\bibitem{ref129}
Zhou et al. (2023). P2P energy sharing and trading in smart communities: pricing, decision-making, agent-based collaboration. \textit{Renewable Energy}. Available online: https://doi.org/10.1016/j.renene.2023.02.125 (accessed on 17 October 2026).
\bibitem{ref130}
Long et al. (2018). P2P energy sharing through two-stage aggregated battery control in community microgrid. \textit{Applied Energy}. Available online: https://doi.org/10.1016/j.apenergy.2018.05.097 (accessed on 17 October 2026).
\bibitem{ref131}
Luo et al. (2022). Distributed P2P energy trading based on game theory, DER ownership complexity. \textit{J. Cleaner Production}. Available online: https://doi.org/10.1016/j.jclepro.2022.131573 (accessed on 17 October 2026).
\bibitem{ref132}
Paudel et al. (2019). P2P Energy Trading in a Prosumer-Based Community Microgrid: Game-Theoretic Model. \textit{IEEE Trans. Industrial Electronics}. Available online: https://doi.org/10.1109/tie.2018.2874578 (accessed on 17 October 2026).
\bibitem{ref133}
Anoh et al. (2020). Energy P2P Trading in Virtual Microgrids: A Game-Theoretic Approach. \textit{IEEE Trans. Smart Grid}. Available online: https://doi.org/10.1109/tsg.2019.2934830 (accessed on 17 October 2026).
\bibitem{ref134}
Yan et al. (2023). P2P transactive energy trading of multiple microgrids under renewable uncertainty. \textit{Int. J. Electrical Power \& Energy Systems}. Available online: https://doi.org/10.1016/j.ijepes.2023.109235 (accessed on 17 October 2026).
\bibitem{ref135}
Alam et al. (2019). Peer-to-peer energy trading among smart homes. \textit{Applied Energy}. Available online: https://doi.org/10.1016/j.apenergy.2019.01.091 (accessed on 17 October 2026).
\bibitem{ref136}
Vieira et al. (2021). P2P Energy Trading in a Microgrid Leveraged by Smart Contracts. \textit{Renewable and Sustainable Energy Reviews}. Available online: https://doi.org/10.1016/j.rser.2021.110900 (accessed on 17 October 2026).
\bibitem{ref137}
Lan et al. (2025). Scalable P2P Energy Trading System With Microgrid Coalitions. \textit{ICSGGE 2025}. Available online: https://doi.org/10.1109/icsgge64667.2025.10985166 (accessed on 17 October 2026).
\bibitem{ref138}
Alfaverh et al. (2023). Dynamic P2P Electricity Market Model for Community Microgrid with Price-Based DR. \textit{IEEE Trans. Smart Grid}. Available online: https://doi.org/10.1109/tsg.2023.3246083 (accessed on 17 October 2026).
\bibitem{ref139}
Wang et al. (2018). P2P Energy Trading Among Microgrids with Multidimensional Willingness. \textit{Energies}. Available online: https://doi.org/10.3390/en11123312 (accessed on 17 October 2026).
\bibitem{ref140}
Lauss et al. (2021). Accurate and Stable HIL Real-Time Simulation of Integrated Power Electronics and Power Systems. \textit{IEEE Trans. Power Electronics}. Available online: https://doi.org/10.1109/tpel.2020.3040071 (accessed on 17 October 2026).
\bibitem{ref141}
Mohammadi et al. (2022). Real-Time Controller-HIL Testing of Power Electronics Converters. \textit{PEDSTC 2022}. Available online: https://doi.org/10.1109/pedstc53976.2022.9767506 (accessed on 17 October 2026).
\bibitem{ref142}
Lu et al. (2007). Low-Cost Real-Time HIL Testing Approach of Power Electronics Controls (VTB-RT). \textit{IEEE Trans. Industrial Electronics}. Available online: https://doi.org/10.1109/tie.2007.892253 (accessed on 17 October 2026).
\bibitem{ref143}
D'Arco et al. (2024). Multi-HIL Laboratory Testing of Power Converters and IEDs for Large-Scale Power Systems. \textit{IEEE Open J. Power Electronics}. Available online: https://doi.org/10.1109/ojpel.2024.3460709 (accessed on 17 October 2026).
\bibitem{ref144}
Bai et al. (2026). Multilayer Device-Level Electro-Thermal Real-Time Simulation and Multipurpose HIL Testing. \textit{IEEE Trans. Power Electronics}. Available online: https://doi.org/10.1109/tpel.2025.3650187 (accessed on 17 October 2026).
\bibitem{ref145}
Lucía et al. (2011). Real-Time FPGA-Based HIL Simulation Test Bench for Multiple-Output Power Converters. \textit{IEEE Trans. Industry Applications}. Available online: https://doi.org/10.1109/tia.2010.2102997 (accessed on 17 October 2026).
\bibitem{ref146}
Estrada et al. (2020). Real-Time HIL Simulation Methodology for Power Converters Using LabVIEW FPGA. \textit{Energies}. Available online: https://doi.org/10.3390/en13020373 (accessed on 17 October 2026).
\bibitem{ref147}
Shekhar et al. (2023). Development of Reliable Power Electronic Systems using Real Time Digital Twin Based PHIL Testbed. \textit{PowerTech 2023}. Available online: https://doi.org/10.1109/powertech55446.2023.10202818 (accessed on 17 October 2026).
\bibitem{ref148}
Badar et al. (2021). Efficient Real-Time Controller Design Test Bench for Power Converter Applications. \textit{IEEE Access}. Available online: https://doi.org/10.1109/access.2021.3107911 (accessed on 17 October 2026).
\bibitem{ref149}
Lauss et al. (2019). Multirate Partitioning Interface for Enhanced Stability of PHIL Real-Time Simulation. \textit{IEEE Trans. Industrial Electronics}. Available online: https://doi.org/10.1109/tie.2018.2826482 (accessed on 17 October 2026).
\bibitem{ref150}
Matar et al. (2010). FPGA Implementation of the Power Electronic Converter Model for Real-Time Simulation. \textit{IEEE Trans. Power Delivery}. Available online: https://doi.org/10.1109/tpwrd.2009.2033603 (accessed on 17 October 2026).
\end{thebibliography}